\documentclass[11pt]{article}

\usepackage[margin=1in]{geometry}
\usepackage{amsmath,amssymb,amsthm}
\usepackage{mathtools}
\usepackage{microtype}
\usepackage{booktabs}
\usepackage{array}
\usepackage{hyperref}
\hypersetup{colorlinks=true,linkcolor=blue,citecolor=blue,urlcolor=blue}

\newcommand{\OHP}{\mathrm{OHP}}
\newcommand{\CS}{C_{\!S}}
\newcommand{\CSeff}{C_{\!S,\mathrm{eff}}}
\newcommand{\dd}{\mathrm{d}}
\newcommand{\eb}{\varepsilon_{b}}
\newcommand{\eS}{\varepsilon_{S}}

\title{The Stern Compact-Layer Robin Boundary Condition\\
       \large Derivation and exact code mapping from
              Bonnefont--Argoul--Bazant (2001)}
\author{PNPInverse / Forward solver}
\date{\today}

\begin{document}
\maketitle

\begin{abstract}
The PNP--BV forward solver imposes a single Robin boundary condition
on the Poisson equation at the electrode,
$\eb\,\partial_{n}\Phi = \CS\,(V_{\mathrm{app}}-\Phi)$, implemented as
the residual line
\begin{verbatim}
F_res -= stern_coeff * (phi_applied - phi) * w * ds(electrode_marker)
\end{verbatim}
in \texttt{Forward/bv\_solver/forms\_logc.py:631} and
\texttt{forms\_logc\_muh.py:668}. This note derives that line, term by
term, from Section~2.2 of Bonnefont, Argoul, and Bazant
(2001)~\cite{bonnefont2001}, the original PNP-with-Faradaic-kinetics
treatment in which this exact boundary condition appears coupled to a
Butler--Volmer flux at the outer Helmholtz plane.
\end{abstract}

%==============================================================
\section{The ``missing'' boundary condition}
%==============================================================

Once the diffuse layer is resolved as a continuum (rather than
collapsed into an effective $\zeta$-potential), Poisson's equation
acquires a degree of freedom at the electrode that the species flux
conditions do not constrain. Bonnefont, Argoul, and Bazant make this
explicit in their Section~2.2~\cite[p.~54]{bonnefont2001}:

\begin{quote}
\emph{``No other boundary conditions are typically mentioned in
electrochemistry textbooks because none are needed when the common
assumption of electroneutrality is made. In this work, however, since
we treat diffuse charge explicitly we also need another boundary
condition on the potential. Ignorance of this boundary condition is
typically hidden in the `zeta-potential' $\zeta$ \dots\ but in
electrochemistry, the zeta potential of an electrode should be
determined self-consistently from the microscopic electrochemical
boundary conditions.''}
\end{quote}

The Stern compact layer supplies this missing condition.

%==============================================================
\section{The Stern boundary condition}
%==============================================================

Bonnefont \emph{et al.}~adopt the Stern model: a thin charge-free
dielectric slab of permittivity $\eS$ and effective thickness $L_{S}$
between the metal and the outer Helmholtz plane, modelled as a
constant areal capacitance
\begin{equation}
\CS \;=\; \eS / L_{S}\qquad [\mathrm{F/m^{2}}].
\label{eq:cs-defn}
\end{equation}
With $\Delta\Phi_{S}\equiv\Phi_{e}-\Phi_{r}$ denoting the potential
drop across the slab (metal minus reaction plane), the paper writes
the boundary condition as~\cite[Eq.~(8)]{bonnefont2001}:
\begin{equation}
\Delta\Phi_{S}\;=\;
\begin{cases}
-\,\eS\,\partial_{X}\Phi / \CS &\text{at } X=0,\\[2pt]
+\,\eS\,\partial_{X}\Phi / \CS &\text{at } X=L,
\end{cases}
\label{eq:bab-eq8}
\end{equation}
together with the definition of the effective Stern length
\begin{equation}
\lambda_{S}\;=\;\eS/\CS,
\label{eq:lambda-s}
\end{equation}
which is~\cite[Eq.~(9)]{bonnefont2001}.

%==============================================================
\section{From \texorpdfstring{Eq.~(\ref{eq:bab-eq8})}{Eq. (8)} to the
standard Robin form}
%==============================================================

Let $\mathbf{n}$ be the outward normal to the PNP domain
$\Omega$ at the electrode boundary $\partial\Omega_{\mathrm{el}}$. At
the right electrode $X=L$, $\mathbf{n}=+\hat{x}$ and
$\partial_{n}\Phi=\partial_{X}\Phi$; at the left electrode $X=0$,
$\mathbf{n}=-\hat{x}$ and $\partial_{n}\Phi=-\partial_{X}\Phi$. Both
cases of Eq.~(\ref{eq:bab-eq8}) then collapse to the single
expression
\begin{equation}
\Delta\Phi_{S}\;=\;\lambda_{S}\,\partial_{n}\Phi.
\label{eq:robin-bab}
\end{equation}
Now $\Delta\Phi_{S}=V_{\mathrm{app}}-\Phi_{\OHP}$, where $\Phi_{\OHP}$
is the boundary value of $\Phi$ on the diffuse-layer side of the
outer Helmholtz plane (the value the PNP solver actually controls).
Multiplying both sides of (\ref{eq:robin-bab}) by $\eb/\lambda_{S}$
gives
\begin{equation}
\boxed{\;\;
\eb\,\partial_{n}\Phi \;=\; \CSeff\,
\bigl(V_{\mathrm{app}} - \Phi_{\OHP}\bigr),
\qquad \CSeff \;\equiv\; \eb/\lambda_{S},
\;\;}
\label{eq:our-robin}
\end{equation}
which is the Robin boundary condition the solver imposes.
Two remarks on the identification of $\CSeff$ with the
physical compact-layer capacitance:

\begin{itemize}
\item In the special case $\eS=\eb$, $\lambda_{S}=L_{S}$ exactly and
$\CSeff=\eb/L_{S}=\CS$ literally.
\item In the general case $\eS\neq\eb$, the displacement field is
continuous across the Stern slab,
$\eb\,\partial_{n}\Phi|_{\mathrm{diffuse}} =
\eS\,\partial_{n}\Phi|_{\mathrm{Stern}}$, so
$\eb\,\partial_{n}\Phi$ on the diffuse side equals
$\sigma_{M}=\CS\,\Delta\Phi_{S}$ (the metal surface charge). The Robin
coefficient appearing in (\ref{eq:our-robin}) is then identified with
the literature value of the compact-layer capacitance directly, and
the role of $\lambda_{S}$ in Eq.~(\ref{eq:bab-eq8}) is to encode the
dielectric mismatch.
\end{itemize}
For the production target $\CS=0.20~\mathrm{F/m^{2}}$
($L_{S}\approx 5~\textup{\AA}$, $\eS\approx 11.3\,\varepsilon_{0}$),
see \texttt{docs/phase6/CMK3\_capacitance\_literature.md}.

%==============================================================
\section{Weak form: term-by-term derivation of the residual line}
%==============================================================

The solver does not implement (\ref{eq:our-robin}) directly; it
substitutes it into the natural-BC slot of the Poisson weak form.
We trace each term of the resulting residual to a specific equation
in~\cite{bonnefont2001}.

Multiply the Poisson equation $-\nabla\!\cdot\!(\eb\,\nabla\Phi)=\rho$
(Bonnefont--Argoul--Bazant Eq.~(3)) by a test function $w$ and
integrate over $\Omega$:
\begin{equation}
-\!\int_{\Omega} \nabla\!\cdot\!(\eb\,\nabla\Phi)\,w\,\dd x
\;=\;\int_{\Omega}\rho\,w\,\dd x.
\label{eq:poisson-weak0}
\end{equation}
Integration by parts (Green's identity) gives
\begin{equation}
\int_{\Omega}\eb\,\nabla\Phi\cdot\nabla w\,\dd x
\;-\;\int_{\partial\Omega}\eb\,(\partial_{n}\Phi)\,w\,\dd s
\;=\;\int_{\Omega}\rho\,w\,\dd x.
\label{eq:poisson-weak-ibp}
\end{equation}
The test space vanishes on the Dirichlet (bulk) boundary, so only the
electrode boundary $\partial\Omega_{\mathrm{el}}$ contributes. On that
boundary we substitute (\ref{eq:our-robin}):
\begin{equation}
\int_{\Omega}\eb\,\nabla\Phi\cdot\nabla w\,\dd x
\;-\;\int_{\partial\Omega_{\mathrm{el}}}
\CS\,(V_{\mathrm{app}}-\Phi)\,w\,\dd s
\;=\;\int_{\Omega}\rho\,w\,\dd x.
\label{eq:poisson-after-robin}
\end{equation}
Moving the source to the left and identifying the assembled residual
$F_{\mathrm{res}}(\Phi,w)$,
\begin{equation}
F_{\mathrm{res}}(\Phi,w)\;=\;
\underbrace{\int_{\Omega}\eb\,\nabla\Phi\cdot\nabla w\,\dd x}_{\textsc{poisson-vol}}
\;\;\underbrace{-\;\int_{\Omega}\rho\,w\,\dd x}_{\textsc{poisson-rhs}}
\;\;\underbrace{-\;\int_{\partial\Omega_{\mathrm{el}}}\CS\,(V_{\mathrm{app}}-\Phi)\,w\,\dd s}_{\textsc{stern}}
\;=\;0.
\label{eq:Fres-final}
\end{equation}

\subsection*{Term-by-term mapping}

Each piece of (\ref{eq:Fres-final}) corresponds to one residual
assignment in the solver and to one equation in~\cite{bonnefont2001}:

\begin{center}
\renewcommand{\arraystretch}{1.35}
\begin{tabular}{>{\raggedright\arraybackslash}p{0.18\linewidth}
                >{\raggedright\arraybackslash}p{0.26\linewidth}
                >{\raggedright\arraybackslash}p{0.30\linewidth}
                >{\raggedright\arraybackslash}p{0.18\linewidth}}
\toprule
\textbf{Term} & \textbf{Math} & \textbf{Code (\texttt{forms\_logc.py})} &
\textbf{Source in \cite{bonnefont2001}}\\
\midrule
\textsc{poisson-vol} &
$\displaystyle\int_{\Omega}\eb\,\nabla\Phi\!\cdot\!\nabla w\,\dd x$ &
\texttt{F\_res += eps\_coeff * dot(grad(phi),} \texttt{grad(w)) * dx}
\;\;(line~600) &
Eq.~(3), volume term after IBP \\
\textsc{poisson-rhs} &
$\displaystyle-\!\int_{\Omega}\rho\,w\,\dd x$ &
\texttt{F\_res -= charge\_rhs * sum(z[i]*ci[i]*w) * dx}
\;\;(line~602) &
Eq.~(3), RHS source $\rho = zF(C_{C}-C_{A})$\\
\textsc{stern} &
$\displaystyle -\!\int_{\partial\Omega_{\mathrm{el}}}\!\!
\CS\,(V_{\mathrm{app}}-\Phi)\,w\,\dd s$ &
\texttt{F\_res -= stern\_coeff *} \texttt{(phi\_applied - phi) * w *} \texttt{ds(electrode\_marker)}
\;\;(line~631) &
Eq.~(8) substituted into the boundary integral of
Eq.~(\ref{eq:poisson-weak-ibp})\\
\bottomrule
\end{tabular}
\end{center}

The Butler--Volmer Faradaic flux (\cite[Eq.~(7)]{bonnefont2001},
applied at the same $\partial\Omega_{\mathrm{el}}$ with
$\eta = \Delta\Phi_{S}-E_{\mathrm{eq}} = V_{\mathrm{app}}-\Phi-E_{\mathrm{eq}}$)
enters the species residuals in the same file at
\texttt{forms\_logc.py:540--590}; this is the second half of the
Bonnefont \emph{et al.}~boundary system and uses the
\emph{same} $\Phi$ that the Stern term controls.

\subsection*{Nondimensional form}

In dimensionless variables $\phi=zF\Phi/RT$, $x=X/L$,
Bonnefont \emph{et al.}~recast the BC as Eq.~(20),
$\Delta\phi_{S} = \mp\,\delta_{D}\,\varepsilon\,\partial_{x}\phi$, with
$\varepsilon=\lambda_{D}/L$ and $\delta_{D}=\lambda_{S}/\lambda_{D}$.
The solver carries the same rescaling internally: see the docstring
in \texttt{Forward/bv\_solver/nondim.py:62--78}, which states
verbatim
\begin{quote}
\texttt{Physical BC: epsilon * grad(phi).n = C\_stern * (phi\_m - phi)}\\
\texttt{stern\_model = C\_stern * V\_T / (F * c\_ref * L)\quad[dimensionless]}
\end{quote}
i.e.\ the dimensionless residual is identical to
(\ref{eq:Fres-final}) with $\eb$, $\rho$, $\CS$, $V_{\mathrm{app}}$
and $\Phi$ all replaced by their nondimensional counterparts.

%==============================================================
\section{Limiting cases}
%==============================================================

Eq.~(\ref{eq:our-robin}) interpolates between two familiar limits:
\begin{itemize}
\item $\CS\to\infty$ (vanishing compact layer): the Robin condition
forces $\Phi\to V_{\mathrm{app}}$ on $\partial\Omega_{\mathrm{el}}$,
which is the legacy no-Stern Dirichlet boundary the solver falls back
to when \texttt{stern\_capacitance\_f\_m2} is \texttt{None}. In
Bonnefont \emph{et al.}~notation this is the Helmholtz limit
$\delta_{D}\to 0$.
\item $\CS\to 0$ (perfectly blocking compact layer): the condition
forces $\eb\,\partial_{n}\Phi\to 0$, i.e.\ no net displacement-current
flux into the metal. In~\cite{bonnefont2001} this is the
``Gouy--Chapman model with $\delta_{D}\to\infty$.''
\end{itemize}
Finite $\CS$ partitions the applied voltage between the compact layer
($\Delta\Phi_{S}=\sigma_{M}/\CS$) and the diffuse layer
($\Phi_{\OHP}$), exactly as motivated in~\cite[\S2.2]{bonnefont2001}.

%==============================================================
\begin{thebibliography}{9}

\bibitem{bonnefont2001}
A.~Bonnefont, F.~Argoul, and M.~Z. Bazant.
\newblock Analysis of diffuse-layer effects on time-dependent interfacial
  kinetics.
\newblock \emph{J.~Electroanal.~Chem.}, 500:52--61, 2001.
\newblock \href{https://doi.org/10.1016/S0022-0728(00)00470-8}%
              {doi:10.1016/S0022-0728(00)00470-8}.

\end{thebibliography}

\end{document}
